\scanpage{74}
% loc: 1–4
\uline{NB} Si $s=\dot x$, \deleted{et si $a$} on définit $T_\lambda=T_\lambda^s$ sur l'ens des $t$ liés à $s$, i.e. des $\dot y$ avec $x\wedge y=e$, par
\[
 T_\lambda(y)=\dot{\overline{y+\lambda x}}
\]

% loc: 5, cancelled
\deleted{(d'où \unsure{$\dot y$} $T_\lambda(y)=(\unsure{-}$}

% loc: 6–9
(Si on remplace \unsure{$a$|$x$} par $-x$, on trouve la $=$ opération $T_\lambda$) En fait, ici \deleted{T} l'opération $T_\lambda$ est induite par \unsure{la|les} transvection de $E$ : \deleted{\unsure{$cx$|$E$} $u\longmapsto{}$} $\mathrm{id}+\varphi^e(x\otimes x)$

% loc: 10–11
Cette structure se transporte sur $C^*(e)$. Il vaut la peine de l'expliciter :

% loc: 12–13
Notons que
\[
 \underbrace{x\overset{e}{\otimes}x}_{=\varphi^e(\deleted{\unsure{x}}x)}
 \circ\underbrace{y\overset{e}{\otimes}y}_{=\varphi^e(y)}
 =y\overset{e}{\otimes}x\,((\unsure{-x|x}\wedge y)/e)
\]

% loc: right of composition
\margin{\uline{NB} $\operatorname{Tr}\,\deleted{\unsure{x}}y\overset{e}{\otimes}x=(\unsure{x|a}\wedge y)/e$}

% loc: 14–15
donc
\[
 \operatorname{Tr}(\deleted{\unsure{x\otimes}}\underset{\deleted{\unsure{x\overset{e}{\otimes}x}}}{u}\circ v)
 =-\bigl((x\wedge y)/e\bigr)^2
\]

% loc: left at 14–19
% slanted
\margin{Que se passerait-il si on choisit sur $C^*$ la structure de graphe définie par ``l'autre'' relation ``non dégénérée pour ttes caractéristiques'' $\operatorname{Tr}uv=+1$ ??}

% loc: 16
donc
\[
 x\wedge y=\pm e\quad\Longleftrightarrow\quad\boxed{\operatorname{Tr}uv=-1}
\]

% loc: 17–19
Les triangles $\{u,v,w\}$ de $C^\times(e)$ sont donc les triples $\{u,v,w\}$ tels que
\[
 \operatorname{Tr}uv=\operatorname{Tr}vw=\operatorname{Tr}wu=-1
\]

% loc: 20–21, bottom edge
Si \deleted{\unsure{on}}
\[
 \begin{aligned}
 x+y+z&=0 & x\wedge y&=y\wedge z=z\wedge x=e\\
 u&=x\overset{e}{\otimes}x,\quad v=y\overset{e}{\otimes}y,\quad & w&=z\overset{e}{\otimes}z=e
 \end{aligned}
\]
% archive sheet number: 73

